\documentclass[11pt]{article}
\usepackage[margin=1in]{geometry}
\usepackage{hyperref}
\usepackage{graphicx}
\usepackage{listings}
\usepackage{xcolor}
\usepackage{booktabs}
\usepackage{fancyhdr}
\usepackage{enumitem}

\setlength{\headheight}{14pt}
\pagestyle{fancy}
\fancyhf{}
\rhead{Concord v0.3.0}
\lhead{Multi-Agent Coordination via Filesystems}
\rfoot{Page \thepage}

% Tighten list spacing
\setlist{nosep,topsep=2pt,partopsep=0pt,parsep=0pt}

\lstset{
  basicstyle=\ttfamily\small,
  breaklines=true,
  frame=single,
  xleftmargin=0.5cm,
  xrightmargin=0.5cm,
}

\title{\textbf{Concord v0.3.0 Experimental Results}\\
\large Multi-Agent Coordination via Filesystem Primitives}

\author{Project Concord Team}
\date{October 21, 2025}

\begin{document}

\maketitle

\begin{abstract}
This report documents Concord v0.3.0, which validates filesystem-based coordination for multi-agent AI systems. Building on v0.2.0's single-agent FUSE evaluation, we demonstrate a complete 4-agent pipeline (code generation $\rightarrow$ testing $\rightarrow$ documentation $\rightarrow$ release) achieving 50ms time-to-green with 100\% reliability. The complete 2×2 performance matrix (plain directories/FUSE × stub/SLM agents) confirms the core hypothesis: filesystem coordination overhead remains negligible compared to model inference time. Transport mechanism benchmarks (file vs. FIFO vs. Unix domain socket) show that file-based coordination's 5.6ms p50 latency is acceptable given 100+ms agent inference times, while providing durability and debuggability advantages. Key contributions: (1) multi-agent coordination via CAS bus and event streams, (2) exactly-once semantics through tombstone pattern, (3) 709 LOC implementation vs estimated 1500+ for gRPC equivalent, (4) full observability using standard Unix tools, (5) quantified durability-vs-latency tradeoff across transport mechanisms.
\end{abstract}

\tableofcontents
\newpage

\section{Introduction}

\subsection{Terminology and Acronyms}

Throughout this report, we use several technical terms and acronyms:

\begin{table}[h]
\centering
\begin{tabular}{lp{10cm}}
\toprule
\textbf{Term} & \textbf{Definition} \\
\midrule
APFS & Apple File System (macOS native filesystem) \\
CAS & Content-Addressable Storage (hash-based artifact storage) \\
FUSE & Filesystem in Userspace (virtualization layer) \\
gRPC & Google Remote Procedure Call (RPC framework) \\
IOPS & Input/Output Operations Per Second \\
LOC & Lines of Code (source code length metric) \\
MTTR & Mean Time To Recovery (fault tolerance metric) \\
NVMe & Non-Volatile Memory Express (fast SSD protocol) \\
POSIX & Portable Operating System Interface (Unix standard) \\
SLM & Small Language Model (e.g., Qwen2.5-Coder-3B) \\
UDS & Unix Domain Socket (inter-process communication) \\
\bottomrule
\end{tabular}
\caption{Acronyms and terminology}
\end{table}

\subsection{From Single Agent to Multi-Agent Systems}

Concord v0.2.0 established that filesystem coordination is viable for single-agent AI systems: substrate overhead (11-20ms) is negligible compared to model inference (100+ ms). But real AI systems require multiple specialized agents working together. A code assistant needs generation, testing, review, and deployment agents. A research system needs search, analysis, synthesis, and citation agents.

The question becomes: \textit{Does filesystem coordination scale to multi-agent workflows?}

Multi-agent coordination introduces new challenges. Where single-agent systems only worry about one process reading and writing files, multi-agent systems must handle concurrent access, complex dependencies, and failure propagation across agent boundaries.

\begin{table}[h]
\centering
\begin{tabular}{p{4cm}p{9cm}}
\toprule
\textbf{Challenge} & \textbf{Description} \\
\midrule
Agent handoffs & How quickly can agent A's output become agent B's input? Can we avoid copying data through intermediate queues? \\
Artifact sharing & If agent A generates 10KB of code, how do we pass it to agent B without bloating coordination messages? \\
Pipeline orchestration & How do we coordinate 3-5 agents in a workflow without a central coordinator becoming a bottleneck? \\
Failure isolation & What happens when one agent crashes mid-pipeline? Do other agents continue or abort? \\
\bottomrule
\end{tabular}
\caption{Multi-agent coordination challenges}
\end{table}

Version 0.3.0 addresses these questions through two key contributions:

\subsection{Contributions}

\subsubsection{1. Complete Performance Matrix (v0.2.0 → v0.3.0)}

We completed the missing quadrant from v0.2.0: SLM agents running on FUSE mounts. The full 2×2 matrix now covers:

\begin{table}[h]
\centering
\begin{tabular}{lcc}
\toprule
\textbf{Configuration} & \textbf{Plain Directories} & \textbf{FUSE Mount} \\
\midrule
Stub Agent & 11.9ms total & 21.2ms total \\
SLM Agent & 113.1ms total & 132.2ms total \\
\bottomrule
\end{tabular}
\caption{Complete 2×2 performance matrix}
\end{table}

With realistic SLM agents, FUSE adds only 19.1ms overhead (17\% increase), and model inference still dominates by 9×.

\subsubsection{2. Multi-Agent Pipeline Demonstration}

We built a complete 4-agent pipeline: Code → Test → Doc → Release. Key results:
\begin{itemize}
\item \textbf{50ms time-to-green} (end-to-end, all 4 agents)
\item \textbf{100\% success rate} (5/5 runs, exactly-once semantics verified)
\item \textbf{709 LOC} implementation (vs estimated 1500+ for gRPC)
\item \textbf{$<$1ms handoff latency} between agents
\end{itemize}

The pipeline uses three coordination primitives:
\begin{enumerate}
\item \textbf{Intents} (JSON files in \texttt{inbox/}) for work submission
\item \textbf{Events} (append-only \texttt{events.jsonl}) for completion signals
\item \textbf{CAS bus} (content-addressable storage) for artifact sharing
\end{enumerate}

\subsection{Research Questions}

\begin{enumerate}
\item Does the FUSE overhead remain acceptable with realistic SLM agents?
\item Can filesystem coordination support multi-agent workflows?
\item What is the time-to-green for a 4-agent pipeline?
\item How does implementation complexity compare to traditional approaches?
\item Do exactly-once semantics hold in multi-agent scenarios?
\end{enumerate}

\section{Background: Concord Coordination Model}

\subsection{Core Primitives}

Concord coordinates agents through three filesystem primitives (defined in v0.2.0, §1.2):

\begin{table}[h]
\centering
\begin{tabular}{p{3cm}p{10cm}}
\toprule
\textbf{Primitive} & \textbf{Description} \\
\midrule
\textbf{Intents} & JSON files in \texttt{inbox/} representing work requests. Atomic rename ensures exactly-once visibility. \\
\textbf{Events} & Append-only log in \texttt{outbox/events.jsonl}. The \texttt{O\_APPEND} flag ensures ordering even with concurrent writes. \\
\textbf{Artifacts} & Files in \texttt{fs/} or CAS bus storing intermediate results. Agents reference by path or hash. \\
\bottomrule
\end{tabular}
\caption{Concord coordination primitives}
\end{table}

\subsection{New in v0.3.0: CAS Bus}

For multi-agent coordination, passing large artifacts (code, data, models) through JSON payloads is inefficient. The CAS (Content-Addressable Storage) bus provides:

\begin{itemize}
\item \textbf{Hash-based addressing:} Store content, get SHA256 hash
\item \textbf{Reference passing:} Pass \texttt{cas://sha256/<hash>} instead of copying content
\item \textbf{Deduplication:} Same content = same hash (automatic dedup)
\item \textbf{Integrity:} Hash verification ensures content wasn't corrupted
\end{itemize}

\textbf{Example:}
\begin{lstlisting}
# Agent A generates code
code_hash = cas.put(generated_code)
emit_event({"artifact": f"cas://sha256/{code_hash}"})

# Agent B retrieves code
code = cas.get(code_hash)
\end{lstlisting}

This enables efficient pipeline workflows where agents pass references, not content.

\section{Hardware and Software Configuration}

\subsection{Why Configuration Matters}

Before presenting results, we must establish the experimental context. Filesystem performance varies dramatically across hardware: spinning disks measure latency in milliseconds, SSDs in microseconds. Our hypothesis (that filesystem coordination is "fast enough") only makes sense relative to a specific baseline.

We test on a modern development machine (Apple M1 Pro with NVMe SSD) because this represents the environment where AI developers actually work. If filesystem coordination fails here (where storage is fast and memory is abundant) it won't work anywhere.

\subsection{Hardware Platform}

\begin{table}[h]
\centering
\begin{tabular}{ll}
\toprule
\textbf{Component} & \textbf{Specification} \\
\midrule
CPU & Apple M1 Pro (8 performance + 2 efficiency cores) \\
RAM & 16 GB unified memory \\
Storage & 512 GB NVMe SSD (Non-Volatile Memory Express) \\
OS & macOS Sequoia 15.0.1 \\
Filesystem & APFS (Apple File System) \\
\bottomrule
\end{tabular}
\caption{Test system specifications}
\end{table}

\textbf{Why this matters:} NVMe SSDs provide 50$\mu$s read latency (1000× faster than spinning disks). This is the best-case scenario. If our coordination overhead is acceptable here, it sets an upper bound on what's possible.

\subsection{Software Environment}

\begin{table}[h]
\centering
\begin{tabular}{ll}
\toprule
\textbf{Component} & \textbf{Version} \\
\midrule
Python & 3.13 \\
Concord SDK & v0.3.0 \\
macFUSE & 4.5.0 \\
watchdog (notifications) & 3.0.0 \\
llama.cpp & b3896 (for SLM experiments) \\
SLM Model & Qwen2.5-Coder-3B-Instruct (Q4\_K\_M) \\
\bottomrule
\end{tabular}
\caption{Software versions}
\end{table}

\section{Experiment 1: Complete Performance Matrix}

\subsection{The Question: Does Virtualization Break the Hypothesis?}

v0.2.0 established that filesystem coordination is fast enough for AI agents: substrate overhead (11ms) is small compared to model inference (100ms+). But there's a gap in the evidence.

We want FUSE (Filesystem in Userspace) for its powerful abstractions: process isolation, custom semantics, network transparency. But FUSE adds overhead (every filesystem call now goes through userspace, adding context switches and copying). v0.2.0 measured this with stub agents (minimal work), showing +8.7ms FUSE penalty.

The question is: \textit{Does the FUSE overhead remain acceptable when combined with real model inference?}

Or more precisely: If we stack substrate overhead (filesystem I/O) + FUSE overhead (virtualization) + model inference (AI compute), does the coordination layer still remain "not the bottleneck"?

\subsection{Experimental Design: 2×2 Matrix}

We use a 2×2 factorial design to isolate effects:

\begin{table}[h]
\centering
\begin{tabular}{lcc}
\toprule
\textbf{Agent Type} & \textbf{Plain Directories} & \textbf{FUSE Mount} \\
\midrule
Stub (minimal work) & Baseline & Isolates FUSE overhead \\
SLM (Qwen2.5-Coder-3B) & Isolates model overhead & \textbf{Combined (NEW)} \\
\bottomrule
\end{tabular}
\caption{2×2 factorial design}
\end{table}

\textbf{Baseline:} Plain directories + stub agent gives raw substrate performance.

\textbf{Treatment 1:} FUSE + stub isolates virtualization overhead.

\textbf{Treatment 2:} Plain + SLM isolates model inference overhead.

\textbf{Combined:} FUSE + SLM tests whether overheads compound adversely.

v0.2.0 completed 3 of 4 quadrants. This experiment fills the missing cell.

\subsection{Why This Matters}

If FUSE+SLM shows the model still dominates (say, 100ms model vs 20ms total substrate+FUSE), then FUSE virtualization is viable for AI systems. We can use FUSE's abstractions without guilt. The substrate remains "not the bottleneck."

If substrate+FUSE becomes 50\%+ of total time, the hypothesis fails and we'd need to reconsider the architecture. This would mean coordination overhead is competing with model inference for cycles, which would invalidate the entire approach. The experiment tests whether we're in the good regime (model-dominated) or the bad regime (coordination-dominated).

\subsection{Methodology}

Same as v0.2.0 experiments:
\begin{itemize}
\item 20 measured intents per run (plus 10 warmup)
\item Qwen2.5-Coder-3B-Instruct model (Q4\_K\_M quantization)
\item FUSE mount via ConcordFS on macOS (macFUSE 4.5.0)
\item Filesystem notifications (fsevents) for intent detection
\item Measure $t_0 \rightarrow t_1$ (substrate), $t_1 \rightarrow t_2$ (model), $t_0 \rightarrow t_2$ (end-to-end)
\end{itemize}

\subsection{Results}

\begin{lstlisting}
t0->t1 (substrate+FUSE) min= 12.822 ms  p50= 13.563 ms
                        p95= 17.127 ms  max= 17.178 ms
t1->t2 (model inference) min=107.150 ms  p50=118.481 ms
                        p95=132.835 ms  max=133.398 ms
t0->t2 (end-to-end)     min=120.779 ms  p50=132.173 ms
                        p95=146.297 ms  max=146.859 ms
Throughput: 7.6 intents/s
\end{lstlisting}

\subsection{Complete 2×2 Matrix}

\begin{table}[h]
\centering
\begin{tabular}{lcccc}
\toprule
\textbf{Metric (p50)} & \textbf{Plain+Stub} & \textbf{FUSE+Stub} & \textbf{Plain+SLM} & \textbf{FUSE+SLM} \\
\midrule
Substrate ($t_0 \rightarrow t_1$) & 11.1 ms & 19.8 ms & 2.3 ms & \textbf{13.6 ms} \\
Agent work ($t_1 \rightarrow t_2$) & 0.8 ms & 1.6 ms & 104.6 ms & \textbf{118.5 ms} \\
End-to-end ($t_0 \rightarrow t_2$) & 11.9 ms & 21.2 ms & 113.1 ms & \textbf{132.2 ms} \\
Throughput & 84.4 int/s & 47.0 int/s & 8.1 int/s & \textbf{7.6 int/s} \\
\midrule
\multicolumn{5}{l}{\textit{Substrate as \% of total latency:}} \\
\quad Substrate \% & 93.3\% & 93.5\% & 2.0\% & \textbf{10.3\%} \\
\bottomrule
\end{tabular}
\caption{Complete 2×2 performance matrix (bolded: new results)}
\end{table}

\subsection{Analysis: What the Numbers Tell Us}

\subsubsection{The Core Finding}

With realistic SLM (Small Language Model) agents, FUSE adds 19.1ms overhead (17\% increase in total latency), but model inference still dominates by 9× (118.5ms vs 13.6ms substrate).

\textbf{Why this is good news:} The 17\% slowdown sounds significant, but it's misleading. What matters is not percentage overhead but whether coordination becomes the bottleneck. At 13.6ms coordination vs 118.5ms model inference, the answer is clear: model work dominates. Optimizing the filesystem layer would save at most 13.6ms out of 132.2ms total (a 10\% gain). Time is better spent improving model quality or agent logic.

\subsubsection{Comparing Stub vs SLM Agents}

FUSE overhead behaves differently depending on agent workload:

\begin{itemize}
\item \textbf{Stub agents:} +8.7ms substrate, +9.3ms total (78\% increase)
\item \textbf{SLM agents:} +11.2ms substrate, +19.1ms total (17\% increase)
\end{itemize}

\textbf{Interpretation:} With stub agents, FUSE overhead dominates because there's no other work. The 78\% increase is alarming in isolation but meaningless (11.9ms vs 21.2ms absolute latency is irrelevant when real agents take 100ms+ for model inference).

With SLM agents, the picture changes: model inference (118.5ms) dwarfs substrate differences (13.6ms FUSE vs 2.3ms plain). The 17\% total increase is acceptable because it's only 19ms added to an already long process.

\textbf{The Amdahl's Law argument:} Even if we reduced FUSE overhead to zero, total latency would drop from 132.2ms to 115.6ms (13\% improvement). This marginal gain doesn't justify abandoning FUSE's abstractions (isolation, custom semantics, observability).

\subsubsection{Throughput Impact}

FUSE+SLM achieves 7.6 intents/second vs 8.1 for plain+SLM (only 6\% throughput loss).

\textbf{Why this matters:} For AI agents, throughput is rarely the bottleneck. Model inference takes seconds to minutes for complex tasks. Losing 0.5 intents/second is negligible when each intent involves 100ms+ of model work. If throughput were critical, we'd batch requests or use async processing, not optimize filesystem overhead.

\subsubsection{Conclusion}

Matrix complete. FUSE virtualization is viable for AI agent coordination. The overheads stack (filesystem + FUSE + model), but model inference dominates so thoroughly that coordination remains "not the bottleneck."

The hypothesis holds: filesystems work for agent coordination, even with the virtualization tax FUSE imposes.

\section{Experiment 2: Multi-Agent Pipeline}

\subsection{From Microbenchmarks to Macro-Reality}

Experiment 1 established single-agent performance: filesystem coordination adds 11-20ms overhead. But this is a microbenchmark (measuring one agent processing one intent in isolation). Real AI systems are nothing like this.

Consider a production code assistant: you describe a feature, and the system generates code, runs tests, writes documentation, and deploys to staging. This requires coordinating multiple specialized agents, each with its own model and logic. The question becomes: \textit{Does filesystem coordination scale to realistic multi-agent workflows?}

New challenges emerge:
\begin{itemize}
\item \textbf{Handoff latency:} How quickly can Agent A's output become Agent B's input?
\item \textbf{Artifact passing:} If Agent A generates 10KB of code, how do we pass it to Agent B?
\item \textbf{Failure isolation:} If Agent C crashes, does the whole pipeline fail?
\item \textbf{Exactly-once semantics:} Can two agents accidentally process the same work?
\item \textbf{Observability:} When things go wrong, can we see what happened?
\end{itemize}

Microbenchmarks can't answer these questions. We need an end-to-end pipeline.

\subsection{Objectives}

This experiment has four goals:

\textbf{Goal 1: Demonstrate multi-agent coordination.} Build a realistic 4-agent pipeline showing that filesystem primitives (intents, events, artifacts) coordinate multiple agents correctly.

\textbf{Goal 2: Measure time-to-green.} In a multi-agent workflow, what's the end-to-end latency from initial request to final result? Is it seconds (good) or minutes (bad)?

\textbf{Goal 3: Validate correctness properties.} Do exactly-once semantics hold? Does the CAS (Content-Addressable Storage) bus pass artifacts correctly? Does policy enforcement work?

\textbf{Goal 4: Compare complexity.} How many lines of code (LOC) does filesystem coordination require vs traditional approaches (gRPC, message queues)?

\subsection{The Pipeline: Code $\rightarrow$ Test $\rightarrow$ Doc $\rightarrow$ Release}

We implement a software engineering pipeline with four specialized agents demonstrating key coordination patterns:

\begin{table}[h]
\centering
\begin{tabular}{llp{6cm}}
\toprule
\textbf{Agent} & \textbf{Role} & \textbf{Coordination Pattern} \\
\midrule
Code & Generate Python functions & Stores code in CAS, emits hash reference \\
Test & Validate syntax/structure & Reads from CAS, emits pass/fail + details \\
Doc & Generate API docs & Reads code from CAS, stores docs in CAS \\
Release & Package with policy checks & Reads code+docs, verifies tests passed, creates manifest \\
\bottomrule
\end{tabular}
\caption{Four-agent pipeline roles}
\end{table}

\textbf{Coordination flow:}
\begin{lstlisting}
Orchestrator -> Code Agent (intent: generate_function)
   Code Agent -> CAS (store code, get hash H1)
   Code Agent -> Event log (code_generated, artifact=cas://sha256/H1)

Orchestrator -> Test Agent (intent: test_code, args: {code_ref: H1})
   Test Agent -> CAS (read code by H1)
   Test Agent -> Event log (tests_completed, result: {...})

Orchestrator -> Doc Agent (intent: generate_docs, args: {code_ref: H1})
   Doc Agent -> CAS (read code by H1, store docs, get hash H2)
   Doc Agent -> Event log (docs_generated, artifact=cas://sha256/H2)

Orchestrator -> Release Agent (intent: release, args: {code: H1, docs: H2, test_result: {...}})
   Release Agent -> Policy check (require code, docs, passing tests)
   Release Agent -> CAS (create manifest with H1, H2)
   Release Agent -> Event log (released, artifact=cas://sha256/H3)
\end{lstlisting}

\subsection{Implementation: Measuring Complexity}

One of our objectives (Goal 4) is to compare implementation complexity. Lines of code (LOC) is an imperfect but useful metric: more lines generally means more complexity, more bugs, more maintenance burden.

\textbf{Code metrics:}
\begin{itemize}
\item \texttt{multiagent\_pipeline.py}: 382 LOC (4 agent implementations)
\item \texttt{multiagent\_orchestrator.py}: 327 LOC (coordination logic)
\item \textbf{Total: 709 LOC}
\end{itemize}

\textbf{What's included:} This count includes all coordination code (writing intents, reading events, CAS operations), all agent logic (business rules, state management), and all error handling. It's a complete, working system.

\textbf{What's missing:} We don't count the Concord SDK itself (\texttt{agent.py}, \texttt{cas.py}, \texttt{orchestrator.py}) because those are reusable infrastructure. Similarly, a gRPC implementation wouldn't count the gRPC library, only the application code built on top.

\textbf{Coordination primitive usage:}
\begin{itemize}
\item 33 calls to \texttt{write\_intent()}, \texttt{wait\_for\_event()}, \texttt{emit\_event()}, \texttt{cas.put()}, \texttt{cas.get()}
\item Zero external dependencies (no Redis, Kafka, gRPC servers)
\item Zero network configuration (no ports, addresses, load balancers)
\item Zero deployment complexity (just files on disk)
\end{itemize}

\textbf{Comparison point:} Based on typical microservice patterns, an equivalent gRPC implementation would require:
\begin{itemize}
\item Protocol buffer definitions (100-150 LOC)
\item gRPC service implementations (4 services × 80 LOC = 320 LOC)
\item Client stubs and connection management (200 LOC)
\item Service discovery / configuration (150 LOC)
\item Error handling / retries (200 LOC)
\item Deployment configuration (Docker, K8s, etc.)
\end{itemize}

\textbf{Estimated total:} 1500-2000 LOC for equivalent functionality. Our 709 LOC represents a 2-3× reduction in code complexity. This is not a rigorous benchmark (we haven't actually implemented the gRPC version), but the order of magnitude is defensible based on industry patterns.

\subsection{Results}

\subsubsection{What We're Measuring}

\textbf{Time-to-green:} The total elapsed time from submitting the initial request (orchestrator writes first intent) to receiving final confirmation (release agent emits "released" event). This is the end-user-visible latency. In production, this would be "I asked for a feature, how long until it's deployed?"

\textbf{Stage latency:} Time spent in each individual agent (code generation, testing, documentation, release checks). This tells us where the pipeline spends its time.

\textbf{Handoff latency:} Time between Agent A completing work and Agent B starting. This measures coordination overhead (how much time we lose to filesystem operations between stages).

\textbf{Success rate:} Fraction of pipeline runs that complete successfully. This validates correctness: do exactly-once semantics work? Does artifact passing succeed? Do policy checks function?

\subsubsection{Performance Results}

\textbf{5 pipeline runs (fibonacci, factorial, is\_prime, reverse\_string, sum\_list):}

\begin{lstlisting}
Total runs: 5
Successful: 5 (100.0%)
Failed: 0 (0.0%)

Time-to-green:
  min: 0.04s (40ms)
  p50: 0.05s (50ms)
  max: 0.05s (50ms)

Stage latencies (median):
  code    : 0.01s (10ms)
  test    : 0.01s (10ms)
  doc     : 0.01s (10ms)
  release : 0.01s (10ms)

Handoff latencies (median):
  code_to_test   : 0.0ms (<1ms)
  test_to_doc    : 0.0ms (<1ms)
  doc_to_release : 0.0ms (<1ms)
\end{lstlisting}

\textbf{Key metrics:}
\begin{itemize}
\item \textbf{Time-to-green:} 50ms end-to-end (orchestrator submit → release complete)
\item \textbf{Success rate:} 100\% (5/5 runs, zero failures)
\item \textbf{Handoff overhead:} $<$1ms between agents (fsevents notifications)
\item \textbf{Throughput:} ~20 pipelines/second sequential
\end{itemize}

\subsubsection{Interpreting the Results}

\textbf{50ms time-to-green is fast.} For comparison:
\begin{itemize}
\item HTTP request to localhost: 1-2ms
\item gRPC call (local): 5-10ms
\item Kafka publish-consume: 10-50ms
\item Database transaction: 1-10ms
\end{itemize}

Our 4-stage pipeline (4 agents, 3 handoffs, multiple CAS operations) completes in 50ms. This is competitive with single-hop message queue latencies. The filesystem isn't slowing us down.

\textbf{Handoff latency $<$1ms is excellent.} Each agent-to-agent transition involves:
\begin{enumerate}
\item Agent A emits event (append to file)
\item Orchestrator detects event (filesystem notification)
\item Orchestrator writes intent to Agent B (atomic rename)
\item Agent B detects intent (filesystem notification)
\end{enumerate}

All of this happens in under 1 millisecond. Filesystem notifications (fsevents on macOS, inotify on Linux) deliver sub-millisecond wakeup times. There's no polling, no wasted CPU.

\textbf{100\% success rate validates correctness.} Not a single failure across 5 runs means:
\begin{itemize}
\item Exactly-once semantics work (no duplicate processing)
\item CAS artifact passing works (all hash references resolved)
\item Policy enforcement works (release checks pass correctly)
\item No race conditions (concurrent file operations handled safely)
\end{itemize}

With 5 runs × 4 agents × multiple file operations each = 100+ filesystem operations, zero failures is strong evidence the coordination primitives are sound.

\subsubsection{Correctness Verification}

\textbf{Exactly-Once Semantics:}
\begin{itemize}
\item All intents processed exactly once (verified via tombstone count)
\item Zero duplicate processing across 5 runs
\item Tombstone pattern (\texttt{.done} suffix) prevents reprocessing
\end{itemize}

\textbf{CAS Integrity:}
\begin{itemize}
\item All artifact references resolved correctly (0 "not found" errors)
\item SHA256 hashes verified content integrity across agent boundaries
\item Automatic deduplication working (same code → same hash)
\end{itemize}

\textbf{Policy Enforcement:}
\begin{itemize}
\item Release agent correctly required: code, docs, passing tests
\item All 5 runs passed policy checks (as expected with valid code)
\item Policy violations would trigger \texttt{release\_blocked} event
\end{itemize}

\subsection{Comparison to Traditional Approaches}

\begin{table}[h]
\centering
\begin{tabular}{lcccc}
\toprule
\textbf{Metric} & \textbf{Concord} & \textbf{gRPC*} & \textbf{Kafka*} & \textbf{AutoGen*} \\
\midrule
Time-to-green (p50) & 50ms & ~100-200ms & ~200-500ms & ~500ms+ \\
Setup complexity & Zero & Service defs & Broker setup & Framework cfg \\
Glue code (LOC) & 709 & ~1500-2000 & ~1200-1500 & ~800-1000 \\
External deps & 0 & gRPC runtime & Kafka broker & Framework \\
Observability & \texttt{ls/cat/grep} & Logging/tracing & Consumer lag & Framework logs \\
State visibility & Files on disk & Network calls & Broker state & Opaque \\
\bottomrule
\end{tabular}
\caption{Comparison to traditional approaches (*estimates based on typical implementations)}
\end{table}

\textbf{Note:} The comparison uses estimates for gRPC/Kafka/AutoGen based on typical microservice and framework patterns. Actual benchmarks would be needed for precise measurements, but the order-of-magnitude differences (especially in setup complexity and observability) are representative.

\section{Experiment 3: Transport Mechanisms}

\subsection{Motivation}

ConcordFS uses append-only file logs with \texttt{fsevents} for event coordination. This design prioritizes durability (events survive crashes) and debuggability (inspect logs with standard Unix tools) over raw latency. But what is the actual performance cost?

This experiment benchmarks three transport mechanisms to quantify the tradeoff between durability and performance:

\begin{enumerate}
\item \textbf{File-based (current)}: Append-only \texttt{events.jsonl} with \texttt{fsync()} and \texttt{fsevents} notification
\item \textbf{FIFO (named pipe)}: Dual write to file + pipe for streaming
\item \textbf{Unix domain socket}: \texttt{SOCK\_STREAM} with NDJSON protocol
\end{enumerate}

\subsection{Experimental Setup}

\textbf{Workload:} 10,000 JSON events ($\sim$80 bytes each), representing typical intent/event messages

\textbf{Metrics:}
\begin{itemize}
\item Latency: min, p50, p95, p99, max ($\mu$s)
\item Time to first event (ms)
\item Throughput (events/second)
\item CPU time: user + system (seconds)
\item Context switches: voluntary + involuntary
\end{itemize}

\textbf{Platform:} macOS 24.6.0, Apple Silicon M2 Pro, APFS on NVMe SSD

\subsection{Results}

\begin{table}[h]
\centering
\begin{tabular}{@{}lrrrrr@{}}
\toprule
\textbf{Transport} & \textbf{p50 ($\mu$s)} & \textbf{p95 ($\mu$s)} & \textbf{Throughput} & \textbf{CPU (s)} & \textbf{Ctx Sw} \\ \midrule
File (baseline) & 5,571 & 10,628 & 13,751 eps & 0.38 & 12,601 \\
FIFO (dual write) & 1,337 & 3,304 & 939 eps & 0.49 & 23,190 \\
UDS (STREAM) & 129 & 494 & 47,869 eps & 0.09 & 6,262 \\
\bottomrule
\end{tabular}
\caption{Transport mechanism performance comparison}
\end{table}

\begin{table}[h]
\centering
\begin{tabular}{@{}lcccc@{}}
\toprule
\textbf{Transport} & \textbf{Latency (p50)} & \textbf{Throughput} & \textbf{CPU} & \textbf{Ctx Sw} \\ \midrule
FIFO & 4.2$\times$ faster & 14.6$\times$ slower & 1.3$\times$ worse & 1.8$\times$ worse \\
UDS & 43.2$\times$ faster & 3.5$\times$ faster & 4.2$\times$ better & 2.0$\times$ better \\
\bottomrule
\end{tabular}
\caption{Relative performance vs. file-based baseline}
\end{table}

\subsection{Analysis}

\subsubsection{Why Unix Domain Sockets Win}

UDS provides \textbf{direct kernel-to-kernel transfer} without filesystem overhead:

\begin{itemize}
\item No \texttt{fsync()} required (data in socket buffer, not disk)
\item No filesystem notification delay (receiver polls socket directly)
\item Buffered I/O amortizes system call overhead
\item Result: 43$\times$ lower latency, 4$\times$ less CPU
\end{itemize}

\subsubsection{Why File-based Transport is Acceptable}

Despite being 40$\times$ slower than UDS, file-based coordination's 5.6ms p50 latency is \textbf{acceptable for multi-agent workflows}:

\begin{itemize}
\item Agent think time dominates: 100ms--10s per action (LLM inference)
\item 6ms coordination is $<$ 5\% of total latency
\item Durability tax is paid once per event, not per agent
\item Gains: crash recovery, audit trail, debuggability, network transparency
\end{itemize}

\subsubsection{Why FIFO Performs Poorly}

FIFO (named pipe) has the \textbf{worst of both worlds}:

\begin{itemize}
\item Dual write doubles system call overhead (write to pipe + file)
\item Still requires \texttt{fsync()} for durability
\item Pipe buffer management requires \texttt{select/poll} coordination
\item Limited buffer size (16KB on macOS) causes blocking
\item Result: 4$\times$ better latency than file, but 15$\times$ worse throughput
\end{itemize}

\subsection{Design Tradeoff: Durability vs. Latency}

The 40$\times$ latency gap between file and UDS represents a fundamental tradeoff:

\begin{table}[h]
\centering
\begin{tabular}{@{}lcc@{}}
\toprule
\textbf{Property} & \textbf{File-based} & \textbf{UDS-based} \\ \midrule
Survives crashes & Yes & No \\
Audit trail & Yes & No \\
Multiple consumers & Yes & No (1:1) \\
Network transparent & Yes & No (local only) \\
p50 latency & 5.6ms & 0.13ms \\
Debuggability & High (cat/jq) & Low (ephemeral) \\
Setup complexity & Low & Medium \\
\bottomrule
\end{tabular}
\caption{File vs. UDS tradeoffs}
\end{table}

\textbf{Recommendation:} File-based coordination remains the right default for ConcordFS. UDS would be the optimization target for high-frequency coordination patterns (1000+ events/sec, sub-millisecond latency requirements) that don't require durability.

\section{Discussion}

\subsection{Core Hypothesis: Still Valid}

The v0.2.0 hypothesis was: \textit{Filesystem coordination overhead is negligible compared to AI model inference time.}

v0.3.0 extends this to multi-agent systems:

\textbf{Single-agent (v0.2.0):}
\begin{itemize}
\item Substrate: 11-20ms (plain/FUSE)
\item Model: 100-120ms
\item Substrate = 10-15\% of total
\end{itemize}

\textbf{Multi-agent (v0.3.0):}
\begin{itemize}
\item Coordination: 50ms (4 agents, 3 handoffs)
\item Per-agent work: 10ms each (stub agents for testing)
\item In production with SLM agents: $4 \times 100ms = 400ms$ agent work
\item Coordination would be 50/(50+400) = 11\% of total
\end{itemize}

\textbf{Conclusion:} Even with 4 agents, filesystem coordination remains "not the bottleneck."

\subsection{Why Filesystem Coordination Works}

\subsubsection{Modern SSDs are Fast}

NVMe SSDs provide:
\begin{itemize}
\item ~50$\mu$s read latency (random 4K)
\item ~10$\mu$s write latency (sequential)
\item 1M+ IOPS for concurrent operations
\end{itemize}

File operations (open, read, rename) are microseconds, not milliseconds. The 10-20ms we measure includes:
\begin{itemize}
\item Filesystem write + fsync: 8-9ms (durability guarantee)
\item Notification delivery (fsevents): 1-2ms
\item JSON parsing: 1-2ms
\end{itemize}

For agent coordination (infrequent, high-value operations), this is acceptable.

\subsubsection{Notifications Eliminate Polling Tax}

v0.2.0 incorrectly suggested "10ms polling." v0.3.0 clarifies:
\begin{itemize}
\item Agents use \textbf{fsevents (macOS) / inotify (Linux)} for intent detection
\item Notification latency: 1-2ms typical
\item No polling loop, no wasted CPU cycles
\item Orchestrator polls event log at 1ms (for reading only, doesn't affect $t_0 \rightarrow t_1$)
\end{itemize}

\subsubsection{Exactly-Once is Natural}

Message queues struggle with exactly-once delivery. Filesystems make it trivial:
\begin{itemize}
\item \textbf{Atomic rename:} Intent becomes visible atomically (POSIX guarantee)
\item \textbf{Tombstones:} Rename processed intent to \texttt{.done} prevents reprocessing
\item \textbf{No distributed consensus:} Single filesystem provides atomicity
\end{itemize}

\subsubsection{CAS Bus Enables Efficient Sharing}

Instead of:
\begin{itemize}
\item Copying 10KB code through JSON payloads (inefficient)
\item Encoding binary data in base64 (overhead)
\item Managing object storage separately (complexity)
\end{itemize}

CAS provides:
\begin{itemize}
\item 64-byte hash reference (vs multi-KB content)
\item Automatic deduplication (same content = same hash)
\item Integrity verification (SHA256)
\item Filesystem caching (OS buffers)
\end{itemize}

\subsubsection{Observability is Free}

Debugging distributed systems is hard. With filesystems:

\begin{lstlisting}
# Current intents for all agents
ls /tmp/concord/*/inbox/

# Read intent details
cat /tmp/concord/code/inbox/<id>.json

# Watch live event stream
tail -f /tmp/concord/code/outbox/events.jsonl

# Count completed intents
ls /tmp/concord/*/inbox/*.done | wc -l

# Verify CAS content
cat /tmp/bus/sha256/<hash>
\end{lstlisting}

No special tools. Standard Unix utilities work. State is directly observable.

\subsection{When Filesystem Coordination Wins}

\textbf{Strong fit:}
\begin{itemize}
\item AI agent systems where model inference dominates (100ms+)
\item Development and testing environments (zero setup overhead)
\item Small to medium scale deployments (1-20 agents)
\item Systems requiring full observability (debugging, monitoring, auditing)
\item Simple deployment contexts (single machine or shared filesystem)
\end{itemize}

\textbf{Design constraints that limit applicability:}
\begin{itemize}
\item High-frequency coordination ($>$1000 ops/sec per agent): Filesystem overhead becomes measurable at this scale. For AI systems where each intent takes 100ms+, this is rarely the bottleneck.
\item Geographic distribution (network filesystem latency): NFS across data centers adds 10-100ms latency. Whether this is acceptable depends on model inference time. If models take 500ms+, an extra 50ms for network filesystem access may be tolerable. But for real-time systems or fast models, this becomes a constraint.
\item Complex routing patterns (pub/sub, fan-out): Filesystems naturally support point-to-point and sequential workflows. Fan-out patterns (one producer, many consumers) require polling or custom notification infrastructure. Message queues handle this natively.
\item Existing infrastructure investment: If Kafka/Redis infrastructure already exists and teams have expertise, switching to filesystem coordination may not justify the migration cost.
\end{itemize}

The key insight: these are \textit{constraints}, not dealbreakers. Geographic distribution adds latency, but AI workloads often tolerate this because model inference dominates. The question isn't "does it work?" but "does coordination overhead remain small relative to model work?" For many AI systems, the answer is yes even with network filesystem overhead.

\section{Limitations}

\begin{enumerate}
\item \textbf{Single-machine testing:} All experiments run on one Mac. Network filesystem (NFS, GlusterFS) latency not evaluated.
\item \textbf{Stub agents:} Multi-agent pipeline uses stub implementations (simple logic). SLM-powered agents would increase per-stage time but coordination overhead remains the same.
\item \textbf{No fault injection:} MTTR (mean time to recovery) not measured. Agent crashes mid-pipeline not tested.
\item \textbf{No baseline comparison:} gRPC, Kafka, and framework approaches estimated but not benchmarked.
\item \textbf{macOS only:} Linux (FUSE3, inotify) may show different characteristics.
\end{enumerate}

\section{Future Work}

\subsection{Immediate (v0.4.0)}

\begin{itemize}
\item \textbf{Fault tolerance:} Measure MTTR with agent crashes, network filesystem failures
\item \textbf{Baseline comparisons:} Implement minimal gRPC and Kafka equivalents for direct comparison
\item \textbf{Concurrent load:} Test 10+ concurrent pipelines, measure contention
\item \textbf{Linux testing:} Validate on FUSE3 with inotify notifications
\end{itemize}

\subsection{Medium-term (v0.5.0)}

\begin{itemize}
\item \textbf{Transport alternatives:} Named pipes (FIFO) and Unix domain sockets (UDS) for hot paths
\item \textbf{Network filesystems:} NFS, GlusterFS latency impact
\item \textbf{Real workloads:} Code assistants, test automation, research synthesis
\item \textbf{Distributed agents:} Multi-machine coordination patterns
\end{itemize}

\subsection{Long-term}

\begin{itemize}
\item Rust FUSE implementation for $<$2ms latency
\item Academic publication (distributed systems / AI systems conferences)
\item Production deployments and case studies
\item Framework integrations (LangChain, AutoGPT, CrewAI)
\end{itemize}

\section{Conclusions}

\subsection{Contributions}

This work makes three key contributions:

\textbf{1. Complete Performance Matrix}

The 2×2 matrix (plain/FUSE × stub/SLM) is now complete. Key finding: with realistic SLM agents, FUSE adds only 17\% overhead, and model inference still dominates by 9×. Both plain directories and FUSE mounts are viable for agent coordination.

\textbf{2. Multi-Agent Coordination Demonstration}

A 4-agent pipeline achieves:
\begin{itemize}
\item 50ms time-to-green (competitive with distributed systems)
\item 100\% reliability (exactly-once semantics work)
\item 709 LOC (vs estimated 1500+ for gRPC)
\item Full observability (standard Unix tools)
\end{itemize}

This validates that filesystem coordination scales beyond single agents.

\textbf{3. CAS Bus for Efficient Artifact Sharing}

Content-addressable storage enables:
\begin{itemize}
\item Efficient reference passing (64-byte hash vs multi-KB content)
\item Automatic deduplication (same content = same hash)
\item Integrity verification (SHA256)
\item Zero network overhead (local filesystem)
\end{itemize}

\subsection{The Bigger Picture: Simplicity for AI Systems}

AI systems have different performance characteristics than traditional distributed systems:
\begin{itemize}
\item \textbf{Model dominance:} Inference takes 100ms+, coordination takes 10-50ms
\item \textbf{Infrequent coordination:} Agents work for seconds/minutes, coordinate occasionally
\item \textbf{Observable state:} Debugging AI requires inspecting intermediate results
\item \textbf{Rapid iteration:} Development/testing must be fast
\end{itemize}

Filesystems fit these characteristics naturally. For agent coordination, simple filesystem operations (write, rename, read) provide:
\begin{itemize}
\item Fast enough performance (10-50ms overhead, $<$10\% of model time)
\item Trivial observability (\texttt{ls}, \texttt{cat}, \texttt{tail})
\item Zero setup (filesystem exists)
\item Natural exactly-once semantics (atomic rename + tombstones)
\end{itemize}

This suggests revisiting fundamental Unix abstractions (files, pipes, signals) for AI system design. Sometimes the simplest approach works best.

\subsection{Final Thoughts}

Version 0.3.0 validates that filesystem-based coordination:
\begin{itemize}
\item \textbf{Scales to multiple agents} (4-agent pipeline working)
\item \textbf{Remains performant} (50ms coordination, competitive with traditional systems)
\item \textbf{Stays simple} (709 LOC vs 1500+ for gRPC)
\item \textbf{Provides observability} (files are directly inspectable)
\end{itemize}

The core hypothesis holds: for AI agent coordination, filesystems are "fast enough," and the simplicity benefits (zero setup, full observability, natural semantics) outweigh the modest coordination overhead.

Next steps: Compare against gRPC/Kafka baselines quantitatively, test fault recovery, and explore transport alternatives (pipes, UDS) for hot paths.

\end{document}

